\chapter{Project Framework, State of the Art, and Needs Analysis}
\minitoc

\section*{Introduction}
\addcontentsline{toc}{section}{Introduction}

This chapter covers where the project came from and what it was asked to fix. It
introduces the institution that hosted the work and the client whose document
the whole thing is built around, then spends most of its length on how that
document actually gets used today.

That last part is worth taking slowly. The requirements listed at the end of
this chapter were not decided in advance. They are what remains once you have
watched the current process fail in specific ways, which is why the failures
come first and the requirements come after. The chapter finishes with the tools
that already exist for this sort of problem, why none of them helps here, and
what was proposed instead.

\section{General Framework of the Project}

This was a final-year project inside the Master's programme in Data Sciences at
Polytech Intl, carried out over a six-month internship. Roughly four of those
months went into building the system. The rest went on framing the subject,
learning the source material well enough to be useful, and writing this report.

Table~\ref{tab:framework} summarises the arrangement.

\begin{table}[h]
\centering
\caption{Framework of the project.}
\label{tab:framework}
\begin{tabular}{|p{0.34\textwidth}|p{0.56\textwidth}|}
\hline
\rowcolor{myblue}\textcolor{white}{\textbf{Item}} & \textcolor{white}{\textbf{Detail}} \\
\hline
Project type & End-of-studies project \\ \hline
Host institution & Polytech Intl, Tunis \\ \hline
Programme & Professional Master's in Data Science \\ \hline
Academic year & 2025 -- 2026 \\ \hline
Duration & Six months \\ \hline
Client organisation & African Development Bank Group \\ \hline
Academic supervisor & Dr.~Marouene Chaieb \\ \hline
Professional supervisors & Mr.~Yasser Ahmad, Mr.~Surya Joymungul \\ \hline
Source document & Delegation of Authority Matrix, over 250 pages \\ \hline
Scope processed & 79-page subset covering three chapters \\ \hline
\end{tabular}
\end{table}

% ---- OPTION A for "what makes it interesting" (pick one, edit freely) ----
The subject sits between two problems that pull in opposite directions.
Extracting a document rewards rules you can check against a printed page: if a
rule is wrong, you can see exactly where and why. Language models offer the
opposite trade, flexibility with no guarantee that what comes out is faithful to
what went in. Most of the engineering in this project came from using both
without letting the second quietly undo the first.

% ---- OPTION B ----
 What kept me on this was how ordinary the document looks and how badly it breaks automated tools. A person reads a page of the DAM without difficulty. Every library I pointed at the same page came back with something plausible and wrong, and finding out why meant learning to read the file the way the printer wrote it rather than the way a reader sees it.

% ---- OPTION C ----
% The interesting part was not building a chatbot. It was working out what a
% chatbot is not allowed to do when the subject is who controls money, and then
% making that restriction hold in code rather than in a prompt.

\section{Host Organisation}

\subsection{Presentation of Polytech Intl}

Polytech Intl, the \'Ecole Polytechnique Internationale Priv\'ee de Tunis, is a
private higher education institution founded in 2013 and based at Rue du Lac
d'Annecy in Les Berges du Lac~1, Tunis. It was set up largely by institutional
figures from the professional world, principally from banking, insurance and
industry, which is visible in how the school approaches training: the intent is
that graduates leave employable, and ideally capable of building something of
their own.

Teaching covers engineering and related disciplines. Programmes run across
computer science, with tracks in big data, business intelligence, software
engineering, IT-finance and embedded systems, alongside mechatronics, industrial
engineering, financial engineering, telecommunications and architecture.
Students are also offered international mobility, with internships and teaching
modules arranged in partnership with schools in Europe, Asia and North America.
The school is a member of the Agence Universitaire de la Francophonie.

% TODO-ADEM: verify the AUF membership and the founding year against the
% school's own site before printing. Both are from public sources rather than
% from the school directly.

The Data Sciences Master's spans statistics and machine learning, data
engineering, and the software side of turning an analysis into something that
runs and keeps running. A final project is meant to show all three. That
expectation is part of why this one was pushed through to a deployed application
instead of stopping at a clean dataset.

\subsection{Project Context and the Client Organisation}

Polytech Intl hosted and supervised the project. The subject came from elsewhere.
The document at the centre of it belongs to the African Development Bank Group,
which acted as the client. There is a certain symmetry in that: a school founded
in part by people from banking, supervising a project built for a bank.

The African Development Bank Group was founded in 1964 and has its headquarters
in Abidjan, C\^ote d'Ivoire. Eighty-one countries are members, 54 of them
African and 27 from outside the continent, and the Bank exists to support
economic development and social progress across its regional member countries.
Table~\ref{tab:afdb} sets out the essentials.

\begin{table}[h]
\centering
\caption{The client organisation at a glance.}
\label{tab:afdb}
\begin{tabular}{|p{0.34\textwidth}|p{0.56\textwidth}|}
\hline
\rowcolor{myblue}\textcolor{white}{\textbf{Item}} & \textcolor{white}{\textbf{Detail}} \\
\hline
Institution & African Development Bank Group \\ \hline
Founded & 1964 \\ \hline
Headquarters & Abidjan, C\^ote d'Ivoire \\ \hline
Membership & 81 countries: 54 regional, 27 non-regional \\ \hline
Purpose & Economic development and social progress across regional member countries \\ \hline
Group entities & African Development Bank, African Development Fund, Nigeria Trust Fund \\ \hline
Regional office concerned & RDGN: Tunisia, Algeria, Morocco, Mauritania, Libya \\ \hline
Regional hub & Tunis \\ \hline
\end{tabular}
\end{table}

% TODO-ADEM: confirm the official name and wording for RDGN with Yasser or
% Surya before this is printed. Also worth one sentence here on how the subject
% reached you -- who proposed it, and what problem they described. Something
% like: "The subject was proposed by [name], who described [the problem as they
% saw it]." One sentence is enough, but it should be the real version.

Working from a real institution's document changed how the project had to be
done. The material could not be tidied up to suit the parser. Correctness meant
whatever the printed pages say, not a score that happened to be easy to compute.
And every claim in this report about what the system pulls out of the document
can be checked against the page it came from. That constraint runs through the
whole of Chapter~2.
